\documentclass[10pt,a4paper]{article}
\usepackage[margin=18mm]{geometry}
\usepackage[T1]{fontenc}
\usepackage[utf8]{inputenc}
\usepackage{helvet}
\renewcommand{\familydefault}{\sfdefault}
\usepackage{xcolor}
\definecolor{accent}{HTML}{254F43}
\usepackage[colorlinks=true,urlcolor=accent]{hyperref}
\usepackage{enumitem}
\usepackage{titlesec}
\pagestyle{empty}
\setlength{\parindent}{0pt}
\setlength{\parskip}{4pt}
\titleformat{\section}{\large\color{accent}}{}{0em}{}[\titlerule]
\titlespacing*{\section}{0pt}{13pt}{6pt}
\setlist[itemize]{leftmargin=14pt,itemsep=2pt,topsep=3pt,parsep=0pt}
\begin{document}
{\Huge Hyogi Kim}\hfill {\small Updated November 2026}\\[5pt]
UNIST $\cdot$ Computer Science and Engineering\\
\href{mailto:jeff4380@unist.ac.kr}{jeff4380@unist.ac.kr} $\cdot$ \href{https://github.com/HYO-AI-UNIST}{GitHub} $\cdot$ \href{https://hyo-ai-unist.github.io}{hyo-ai-unist.github.io}
\section{Research interests}
AI memory $\cdot$ Continual learning $\cdot$ Robot learning\\
I study computer science at UNIST. My interests lie in AI memory, continual learning, and abstraction: how an agent can preserve useful experience and draw on it when learning something new.
\section{Education}
\textbf{UNIST}\\
Computer Science and Engineering
\section{Selected research \& project experience}
\textbf{Continual learning for VLA}\hfill {\small Ongoing research}\\
{\small Research participant; research questions and method exploration.}
\begin{itemize}
\item Participating in research on sequential VLA learning using LIBERO and SO-101.
\item Exploring experience reuse, capability retention, and the storage and training costs of replay.
\end{itemize}
\vspace{3pt}
\textbf{Maestro / MobileKT}\hfill {\small Project}\\
{\small Team lead $\cdot$ MobileKT model architecture $\cdot$ on-device deployment.}
\begin{itemize}
\item Team lead; responsible for MobileKT model architecture and on-device deployment.
\item Team implementation: 1.48 MB of bundled ONNX models; median knowledge-tracing step latency of 49 ms over 63 steps on a Galaxy Tab S8.
\end{itemize}
\vspace{3pt}
\textbf{Edge AI video processing}\hfill {\small Engineering project}\\
{\small Model and inference-pipeline development.}
\begin{itemize}
\item Developed object-detection and video-processing components for RK3588 and DEEPX DX-M1.
\item Worked with C++ asynchronous inference, model conversion, preprocessing, postprocessing, and video input/output.
\end{itemize}
\vspace{3pt}
\textbf{Nano-MemGPT}\hfill {\small In development}\\
{\small Research proposal and experimental scaffold.}
\begin{itemize}
\item Proposed diagnosis of small-model memory-management failures using oracle function calls and failure categories.
\item Developing an experimental scaffold for teacher trajectories and full fine-tuning / LoRA comparisons.
\end{itemize}
\vspace{3pt}
\section{Technical experience}
\textbf{Model development:} Python $\cdot$ PyTorch $\cdot$ model training and evaluation\\
\textbf{Deployment \& systems:} C++ $\cdot$ ONNX $\cdot$ Linux $\cdot$ asynchronous inference pipelines\\
\textbf{Research environments:} LIBERO $\cdot$ SO-101 $\cdot$ mobile and edge devices\\
\end{document}
